\section{Discussion and Limitations}
\label{sec:discussion}

QuaRK is designed to make resource choices visible rather than to argue that
one quantum reservoir family dominates classical alternatives.  The
architecture $\arch$ separates representation width, reservoir multiplicity,
and memory range, while $M$, $S$, and $\Lambda$ govern measurement precision,
frozen-design search, and readout capacity.  The final risk theorem reflects
that separation: width controls an information floor, the frozen pool controls
the chance of reaching near that floor, and the remaining resources control
statistical, measurement, and finite-window discrepancies.

The projection-information floor is target dependent.  Under the canonical
Gaussian coupling it cannot increase with width, and for $n\geq d$ it vanishes
almost surely because the projection is injective.  This is not a claim that
each additional qubit strictly improves finite-sample prediction.  A target
may already factor through a narrower projection.  Larger $n$ still increases
physical circuit width and the number of exposed observables, but the RMS
Mat\'ern metric prevents that larger coordinate count from mechanically
inflating kernel distances or the shot/window risk-sensitivity bounds.  The
empirical width study is therefore intended to identify tasks for which the
additional projected information is actually useful.

The finite-data projection theorem answers a different question.
Johnson--Lindenstrauss preservation concerns only the finite raw inputs that
appear in the observed windows and can permit $n\ll d$.  It does not eliminate
the population information floor and is not used in the end-to-end risk proof.
Conversely, the population result does not prescribe the smallest useful width
for a particular finite dataset.  Keeping these statements separate avoids
turning finite-set geometry preservation into an unsupported performance
claim.

The random-pool result also has a precise limitation.  Once an architecture
contains a strictly good support point, continuity and full support imply a
positive target-dependent mass $\hitMass$, and the probability of
missing every good design is exactly $(1-\hitMass)^S$.  The paper does
not provide a distribution-free numerical lower bound on
$\hitMass$.  Such a bound would require quantitative small-ball
control for the full random circuit and memory-rate law around target-useful
parameter regions and could be extremely conservative.  The theorem therefore
explains the role and asymptotic leverage of $S$ without pretending that a
universal finite pool size can be prescribed for every task.

The mixer theorem should likewise be read at the correct level.  The identity
support point remains sufficient for the asymptotic Stone--Weierstrass
construction once multiplexing and a universal readout are unrestricted.
What the nontrivial $W_n$ adds is a finite-resource mechanism: generic edge
gates spread seed observables to larger Pauli supports, while the cycle bounds
that spreading by graph neighborhoods and the edge bank supplies two-site
seeds.  This does not imply that entanglement monotonically improves predictive
risk, nor that the cycle is an expressivity-optimal topology.  The mixer
ablation therefore reports operator spreading and feature-rank diagnostics
alongside task error; if spreading increases without a predictive benefit, the
architectural hypothesis is rejected for that regime.

The statistical and physical fidelity bounds are sufficient and can be
conservative.  The finite-shot term uses coordinatewise variance bounds, the
dependent-data term uses independent-block coupling and a uniform RKHS bound,
and the finite-window term uses worst-case contraction.  Their value is
interpretability rather than numerical tightness.  In particular, the
condition $\delta_g>0$ can require the raw-time gap to grow with sample size
when $\beta_{\processZ}(g)$ is not negligible.  This restriction is explicit
rather than asymptotically hidden: under the geometric envelope
$\beta_{\processZ}(q)\leq\mixConst e^{-\mixDecay q}$,
\cref{cor:geometric-mixing-gap} shows that $g=O(\log N)$ suffices at fixed
confidence.  The constants remain process dependent, so the empirical study
reports observed counterparts of these gaps rather than treating the
worst-case bounds as fitted predictions.

Finally, the core empirical program remains inside a controlled beta-mixing,
noiseless task family, where the Bayes risk is known and resource--difficulty
interactions can be isolated.  Good behavior on chaotic or noisy tasks would
be a robustness result, not an extension of the theorem.  The resource
accounting uses the number of independent state preparations/circuit runs as
the acquisition metric under a reusable-$n$-qubit hardware convention.  This
should not be confused with the resources of one run: a length-$w$ trajectory
has logical depth $O(w)$ and width $O(n)$.  Fully parallelizing the $R$ branches
would instead trade wall-clock executions for $O(Rn)$ width.  Hardware-specific
gate times, calibration drift, and
fault-tolerant overhead remain outside scope.  No claim of quantum advantage
follows from the current theory alone.
